\subsection{Objetivos del Proyecto}

Para guiar el ciclo de desarrollo del software y evaluar el cumplimiento de las metas planteadas, se formulan los objetivos generales y específicos que orientan a Mi Trufi.

\subsubsection{Objetivo General}

Diseñar e implementar una plataforma tecnológica multiplataforma integrada por una aplicación web administrativa y una aplicación móvil operativa, orientada a optimizar la gestión de turnos de salida cada 15 minutos, el monitoreo satelital en tiempo real de las unidades y la solicitud anticipada de paradas en el corredor de transporte interurbano Calamuchita -- Tarija.

\subsubsection{Objetivos Específicos}

\begin{enumerate}[label=\arabic*.]
    \item \textbf{Desarrollar el panel web de control y secretaría:} Construir una interfaz de escritorio que permita registrar la flota de vehículos, programar la rotación de turnos a intervalos de 15 minutos y monitorear visualmente la totalidad de unidades activas sobre el mapa georreferenciado de la ruta.
    \item \textbf{Implementar la aplicación móvil para conductores:} Diseñar una aplicación móvil ergonómica y de bajo consumo de recursos que emita las coordenadas GPS del vehículo en tiempo real y permita al chofer actualizar con facilidad su nivel de ocupación (disponible, cupo parcial o lleno).
    \item \textbf{Implementar la aplicación móvil para pasajeros:} Desarrollar una interfaz intuitiva para dispositivos móviles mediante la cual los usuarios puedan rastrear la aproximación de las unidades en trayecto y activar la funcionalidad de <<levantar la mano virtual>> para marcar su parada.
    \item \textbf{Diseñar el motor de geolocalización y notificaciones en tiempo real:} Establecer una arquitectura de comunicación bidireccional basada en sockets o mensajería en tiempo real que transmita eventos de parada solicitada y coordenadas con latencia mínima.
    \item \textbf{Validar el sistema en condiciones operativas reales:} Evaluar la estabilidad del software, la precisión del rastreo y la usabilidad de las aplicaciones en un entorno piloto sobre la carretera entre Calamuchita y Tarija, midiendo la reducción de tiempos muertos y la satisfacción de los actores.
\end{enumerate}
